\documentclass[12pt]{article}
\usepackage[utf8]{inputenc}
\usepackage[T1]{fontenc}
\usepackage{hyperref}
\title{Ulwa Language Resources for K-12 Curriculum Development}
\author{Compiled by Claude Sonnet 5.5 (claude-sonnet-5-5) for Andr\'as Kornai}
\date{}
\begin{document}
\maketitle

\section{Introduction}
Ulwa (ISO 639-3: yla; also called Yaul; Glottocode yaul1241) is a severely endangered Keram language of the East Sepik Province, Papua New Guinea. It is spoken in four villages of the Angoram District: Manu, Maruat, Dimiri and Yaul. Estimates of the number of speakers vary: Barlow's grammar describes about 600 speakers, Wikipedia gives fewer than 700 fluent and about 1,250 semi-fluent, and a rough 2015 estimate suggested about 300 fluent. Children are not currently learning the language. The Manu variety differs considerably from the other three. The search found no evidence of school teaching.

\section{Primary Educational Resources}
The search found no primers, textbooks or school materials.

\section{Linguistic and Reference Materials}
\begin{itemize}
\item Russell Barlow, \emph{A grammar of Ulwa (Papua New Guinea)}, Language Science Press, which includes a lexicon of over 1,400 entries: \url{https://langsci-press.org/catalog/book/372}
\item Wikipedia article on the Yaul (Ulwa) language: \url{https://en.wikipedia.org/wiki/Yaul_language}
\item Refubium (FU Berlin) repository record, which appeared in search results for Ulwa. Its exact content was not verified: \url{https://refubium.fu-berlin.de/handle/fub188/41180?show=full}
\end{itemize}

\section{Religious and Community Resources}
\begin{itemize}
\item Joshua Project profile of the Yaul in Papua New Guinea: \url{https://joshuaproject.net/people_groups/16004}
\end{itemize}

\section{Digital and Audio Resources}
\begin{itemize}
\item Endangered Languages Project entry: \url{https://endangeredlanguages.com/lang/10509}
\item Russell Barlow's staff page at the Max Planck Institute for Evolutionary Anthropology (research on Keram languages): \url{https://www.eva.mpg.de/linguistic-and-cultural-evolution/staff/russell-barlow/}
\end{itemize}
The search did not show whether the recorded texts behind the grammar are archived and accessible.

\section{Curriculum Development Context}
Barlow's grammar and lexicon are a substantial base for any orthography or teaching materials. No educational materials were found. With few young speakers, revitalisation work such as documentation, an orthography and community-led teaching would need to come before a full K-12 curriculum. The Manu dialect differences should be taken into account.

\end{document}